
\(T_i\) sur \(f_i^{(1)}\). Cet
homomorphisme définit une immersion fermée \(X'\longrightarrow\bfP_X^{d-1}\), régulière d'après 1.8 ii).

D'après 1.8 iii), \(\underline{S}\) est parfait en tant que \(\underline{O}_X[T_1,\ldots,T_d]\)-Module gradué. En particulier, \(\underline{S}\) est de présentation finie comme Module gradué. Soit donc:
\[
\underline{L}''\longrightarrow \underline{L}'\longrightarrow \underline{S}\longrightarrow0
\]
une présentation finie graduée de \(\underline{S}\) sur \(\underline{O}_X[T_1,\ldots,T_d]\). On obtient sur \(\bfP_X^{d-1}\) la suite exacte:
\[
\underline{M}''\longrightarrow \underline{M}'\longrightarrow\underline{O}_{X'}\longrightarrow0,
\]
où \(\underline{M}''=\underline{\Proj}(\underline{L}'')\), \(\underline{M}'=\underline{\Proj}(\underline{L}')\), et sont donc localement libres de type fini sur \(\bfP_X^{d-1}\). On peut donc écrire pour tout \(n\) le diagramme commutatif:
\[
\begin{tikzcd}[column sep=large,row sep=large]
\underline{L}''_n \arrow[r] \arrow[d,"\sim"'] &
\underline{L}'_n \arrow[r] \arrow[d,"\sim"'] &
\underline{J}^n \arrow[r] \arrow[d] &
0\\
p_*(\underline{M}''(n)) \arrow[r] &
p_*(\underline{M}'(n)) \arrow[r] &
p_*(\underline{O}_{X'}(n)) \arrow[r] &
0,
\end{tikzcd}
\]
où \(p\) est le morphisme structural: \(\bfP_X^{d-1}\to X\), et \(\underline{L}''_n\), \(\underline{L}'_n\) les composantes de degré \(n\) de \(\underline{L}''\), \(\underline{L}'\). La ligne supérieure est exacte par hypothèse; les deux premières flèches verticales sont des isomorphismes à cause des résultats classiques pour les images directes dans \(\underline{O}_P(n)\), \(P\) étant un fibré projectif (EGA III, 2.1.15); enfin, la dernière ligne est exacte si \(n\) est assez grand, car le noyau de \(\underline{M}'\to\underline{O}_{X'}\) est un \(\underline{O}_P\)-Module parfait, et on peut appliquer le lemme rappelé plus haut, qui montre que ce noyau, tensorisé par \(\underline{O}_P(n)\) pour \(n\) assez grand, est \(p_*\)-acyclique; de même pour le noyau de \(\underline{M}''\to \underline{M}'\). Il en résulte alors que pour \(n\) assez grand, \(\underline{J}^n\longrightarrow p_*(\underline{O}_{X'}(n))\) est un isomorphisme. Or \(p_*(\underline{O}_{X'}(n))=f_*(\underline{O}_{X'}(n))\); on a
